\scanpage{53}
\mspage{26}
3. \uline{Étude des sections de} $\mathcal{E}_U$ \uline{sur} $\Gamma$.

Soit $U$ un \deleted{modèle élémentaire sur} \inserted{schéma connexe lisse de t.f. \inserted{\unsure{géom. conn|géom. conn.|géom. connexe|géométriquem.}}} le corps $K$, d'où $\mathcal{E}_U-\mathcal{E}_K$, et on se propose d'étudier les sections mod $\unsure{\Pi_{U/K}|\pi_{U/K}}$-conjugaison --- plus généralement, on veut en $=$ temps regarder les sections $\mathcal{E}'_K\longrightarrow\mathcal{E}_U$, où $\mathcal{E}'_K$ est un sous-groupe ouvert de $\mathcal{E}_U$ (ce qui signifie que l'on fait une extension de base finie sur $K$). Si \inserted{$K$ de type fini sur le corps $\mathbb{Q}$ et si} $U$ se plonge dans un schéma \deleted{abélien} \inserted{ou un groupe commutatif \deleted{\unsure{algébrique}} \unsure{rigide|rigides}} l'application
\[
U(K)\longrightarrow \inserted{\text{Classes d'isomorphie}}\ \text{Sections de }B_U\text{ sur }B_K
\]
i.e. classes de $\unsure{\Pi_{U/K}|\pi_{U/K}}$-conjugaison \unsure{des|de} sections de $\unsure{B_U|\mathcal{E}_U}$ sur $\unsure{B_K|\mathcal{E}_K}$

est injective. On va examiner d'autres façons \inserted{``géométriques''} de trouver des sections.

Supposons d'abord que $U$ est une courbe algébrique, \deleted{\unsure{que je suppose de|qui je suppose de}} \inserted{qui ne soit pas de} type $(0,0)$ ou $(0,1)$ i.e. $\pi_1(\bar U)=\unsure{\Pi_{U/K}|\pi_{U/K}}\ne\{0\}$ ; Alors \unsure{on a} que pour $\unsure{\forall|A}\ i\in\widehat U\setminus\bar U$ (``pt à l'infini''), le groupe de lacets $L_i$ fournit un scindage (au moins si $i$ est $/K$ i.e fixé

\margin{On a choisi un rev. universel $\widetilde U$ de $U$ pour définir $\bar K$, et $\mathcal{E}_U,\mathcal{E}_K$, et $\mathcal{E}_U-\mathcal{E}_K$}
\margin{\diagram{\begin{tikzcd}[ampersand replacement=\&,row sep=small,column sep=small]
1\&\\
\mathcal{E}_K\arrow[u]\arrow[r,equal]\&\Gamma\\
\mathcal{E}_U\arrow[u,no head]\&\\
\unsure{\Pi_{U/K}|\pi_{U/K}}\arrow[u]\&\\
1\arrow[u]\&\\
(1)\&
\end{tikzcd}}}
\margin{52}
